\subsection{Moving a neighbor operator through the period}
\label{subsec:hecke-period}

The column lattice is now fixed. To use its local divisibility inside the period, one must commute the neighbor operators with the kernels at the stages assigned to their orbits.

\newcommand{\hecbin}[2]{\genfrac{[}{]}{0pt}{}{#1}{#2}_{Q}}

A column generator in \eqref{eq:tag-c4} contains \(m\) copies of the same
operator \(Y_v\) for each orbit \(v\) in a subset \(J\) of the
selected orbits. The stage at which to use a copy can be different
at the different orbits. We derive the operator identity and the
cost of making that choice. The local kernels in this calculation
retain the test function in the theta series.

Fix one of the split general linear factors in the unramified
lattice model. Its residue cardinality is \(Q=p^{2h}\), where the
length of the embedding cycle is \(2h\). Number the stages by
\(1\le j\le m\), put \(g=t+1\) and 
\[
  N_j=k+jg,
\]
 and write
\(K_j\) for the lift from rank \(N_{j-1}\) to rank \(N_j\).
In this paragraph write \(N=N_{j-1}\), \(N'=N_j\).
At the lower and upper lattices \(A,B\) let the test function of
the split place be the indicator of \(zA\subset B\) on
\({\rm Hom}(F^N,F^{N'})\). Write \(l(A),l(B)\) for log volumes
to base \(Q\), relative to the reference lattices. In standard
split coordinates the scalar in the kernel has local part
\[
 a^{l(A)}(a')^{l(B)}
 p^{-h(Nl(B)-N'l(A))}
 p^{\Delta_j(l(A)-l(B))},\qquad
 \Delta_j=\sum_{s\ {\rm even}}w_j(s)-\sum_{s\ {\rm odd}}w_j(s).
 \tag{P1}\label{eq:tag-p1}
\]
\begin{proposition}[The neighbor identity on a single lift]\label{prop:hecke-period-neighbor-identity}
At the chosen split factor with \(Q=p^{2h}\), the source and target neighbor operators satisfy \textup{\eqref{eq:tag-p3}} with the test function retained in \(K_j\),
\[
 E_{j,i}K_j=u_j^iK_j\sum_{l=0}^i e_l\big((m_j\beta_b)_{b=0}^t\big)E_{j-1,i-l}.
\]

Both factors lie away from the character support, the
translation level is prime to \(p\), and the other local
data are fixed. The transfer uses
the same mass of each incident pair in both orders of summation.
\end{proposition}

\begin{proof}
\proofstep{The scalar of the incident vacuum.}

The factors \(a,a'\) are the values on the uniformizer of the
unramified character factors at this stage, with valuation zero.
Only the indicated place varies. More precisely, the kernel on
the indicator includes the fixed linear functionals at the other places, the differential operator and the sum over rational
arguments. The Haar volume of \({\rm Hom}(A,B)\) relative to the
reference ball is \(Q^{Nl(B)-N'l(A)}\), by changing a basis on
each side.

Taking the inverse square root of the Hom-ball volume gives the factor
\[
  \big(Q^{N l(B)-N'l(A)}\big)^{-1/2}=Q^{(N'l(A)-N l(B))/2}.
\]
 The
character of a change of basis with valuation-zero determinant is
trivial here, so the character factors depend only on the two
lengths. The last factor is the determinant twist in the source and target evaluations. In the determinant twist each increment of the length difference has the same multiplier. With \(u=l(A)-l(B)\) the ratio is
\[
  p^{\Delta_j(u+1)}/p^{\Delta_j u}=p^{\Delta_j},\qquad \Delta_j=\sum_{s\ {\rm even}}w_j(s)-\sum_{s\ {\rm odd}}w_j(s).
\]
 These are the hypotheses on the split vacuum
and its normalization in \eqref{eq:tag-p1}.

\proofstep{Apply the coefficientwise indicator relation.}

Proposition~\ref{prop:incident-covariance} supplies the linearity.
Fix the away data on a finite family of lattice pairs, and
use a single pair as reference. At the split prime suppose both members and their conjugates
are away from the character support and the translation level
is prime to the residue characteristic. Denote by \(I_{A,B}\)
the Hom-ball indicator and by \(c_{A,B}\) the full scalar of
\eqref{eq:tag-p1}, including the relative volume to power \(-1/2\).
The source and target lattice changes commute on the split Hom
space.

The kernel \(K_j(A,B)\) for the fixed away data is thus determined by one functional \(\mathcal L\) on the indicators
\[
  c_{A,B}^{-1}K_j(A,B)=\mathcal L(I_{A,B}).
\]
Torsion translation, evaluation and jets all enter this common functional. Thus, for the finite set of pairs in \eqref{eq:tag-p2},
\[
 \sum_{A,B}n_{A,B}I_{A,B}=0
 \quad\Longrightarrow\quad
 \sum_{A,B}n_{A,B}\mathcal L(I_{A,B})=0 .
\]
The implication for rational coefficients also descends to any
algebraic coefficient field in which the values are read.
The identification of the local coordinate lattices requires
that their completions in the two fixed frames equal the lattices
of the Hom model.

\proofstep{Count neighbors by the kernel of the residue map.}

Let \(T_{j,i}\) sum over the sublattices of a target lattice of
codimension \(i\) in its reduction modulo \(p\); \(T_{j,0}=1\).
For fixed \(z\) the number of such targets at rank \(N'\)
containing \(zA\) is
\[
 \sum_{l=0}^i Q^{(i-l)(g-l)}
 \hecbin{g}{l}\,
 \#\{A_1\supset A:[A_1:A]=Q^{i-l},\ pA_1\subset A,                                      zA_1\subset B\},
 \tag{P2}\label{eq:tag-p2}
\]
where terms with \(l>g\) are zero. Indeed if \(zA\not\subset B\),
both sides vanish. Otherwise let \(d\) be the dimension of the
kernel of \(\bar z:A/pA\to B/pB\). The count on the left is
\(\hecbin{d+g}{i}\), by duality in the cokernel. The inner
count on the right is \(\hecbin{d}{i-l}\).

To count dimension
\(i\) planes in a sum of dimensions \(d\) and \(g\), first fix
an \(l\)-plane in the second factor and an \((i-l)\)-plane in
the first. The planes with that intersection and projection are
graphs of maps from the fixed projection in the first factor to
the second factor modulo the fixed intersection, of dimension \(g-l\). This gives
the power of \(Q\) in \eqref{eq:tag-p2} for every rank of \(\bar z\).

\proofstep{Transfer the identity with the same pair masses.}

Shrinking \(B\) by \(i\) multiplies \eqref{eq:tag-p1} by

\[
  C_A=(a')^{-i}p^{i\Delta_j+hNi};
\]
 enlarging \(A\) by
\(i-l\) multiplies it by
\(C_C=a^{i-l}p^{(i-l)\Delta_j+hN'(i-l)}\).
Apply the same linear functionals to \eqref{eq:tag-p2}, multiply by \(C_A\),
and divide the term indexed by \(l\) by \(C_C\).
In the sum over source classes the up correspondence moves to a
down correspondence on the input section. Indeed the mass of
each orbit of incident pairs is its inverse stabilizer order in
either summation, by the orbit-stabilizer formula. Thus all the
class masses in the new sum agree with the original ones. With
\[
 E_{j,i}=p^{-hi(N_j-i)}T_{j,i},\quad
 u_j=(aa')^{-1},\quad m_j=a p^{\Delta_j},\quad
 \beta_b=p^{h(t-2b)}\quad(0\le b\le t)
\]
write \(e_l\) for the elementary symmetric polynomial of degree
\(l\). The identity becomes
\[
 E_{j,i}K_j
 =u_j^i K_j\sum_{l=0}^i
       e_l\big((m_j\beta_b)_{b=0}^t\big) E_{j-1,i-l}.
 \tag{P3}\label{eq:tag-p3}
\]

To check the scalar, the
coefficient of \(z^l\) in \(\prod_{b=0}^t(1+\beta_b z)\)
is \(p^{-h l(g-l)}\hecbin{g}{l}\), by expanding the product
in increasing powers of \(Q\). Multiplying the two sides of
\eqref{eq:tag-p2} by their length and \(E\) normalizations gives exactly
this coefficient times \(u_j^i m_j^l\) on the normalized
source operator. Thus \eqref{eq:tag-p3} includes all ranks, also the identity
operator and the zero coefficient when \(l>g\).

\end{proof}

\paragraph{Removing the extra Euler factors.}

Moving a neighbor operator along the chain introduces Euler factors. Their inverse must act integrally on the same stages for the divisibility gain to survive.

\begin{proposition}[The integral Euler inverse on the chain]\label{prop:hecke-period-euler-inverse}
For \(1\le r\le k/2\), take dominant weights \(\lambda_s\) at the
\(2h\) letters. Assume \(w_1(s)\ge-\lambda_{s,k}\) and
\(w_1(s)\le\cdots\le w_m(s)\). At stage \(1\le j\le m\), the
normalized operator in \textup{\eqref{eq:tag-p5}} transfers \(Y_v\) through
\(K_{[j]}\) with denominator exponent \(P_{j,v}\). For the coefficients \(c_l\), with \(c_r\) the coefficient of the zero-index term,
\[
 v_p(c_l)\ge0\quad(0\le l\le r),\qquad v_p(c_r)\ge A_+.
\]

\end{proposition}

\begin{proof}
\proofstep{Invert the scalar polynomial through degree \(r\).}

Let \(K_{[j]}=K_j\cdots K_1\), \(U_j=u_1\cdots u_j\), \(U_0=1\).
Put \(\mathcal E_j(z)=\sum_{i\ge0}U_j^{-i}E_{j,i}z^i\).
After dividing \eqref{eq:tag-p3} by \(U_j^i\), iteration gives
\[
 \mathcal E_j(z)K_{[j]}=
 K_{[j]}\mathcal E_0(z)
 \prod_{a=1}^j\prod_{b=0}^t(1+\xi_{ab}z),\qquad
 \xi_{ab}=U_{a-1}^{-1}m_a\beta_b .
 \tag{P4}\label{eq:tag-p4}
\]
All products can be taken modulo \(z^{r+1}\), where \(1\le r\le k/2\)
is the neighbor dimension. The inverse of the scalar product has
coefficient \((-1)^l h_l(\xi)\) in degree \(l\), with \(h_l\)
the complete homogeneous symmetric polynomial. This follows by
multiplying the geometric series of its individual factors.

\proofstep{Compute the exponent of \(p\) in the operator identity.}

Use the dominant base weights \(\lambda_{s,1},\ldots,\lambda_{s,k}\)
at the \(2h\) letters and put
\[
 d_i(j)=\sum_{s\ {\rm even}}\sum_{a=1}^i\lambda_{s,a}
             +i\sum_{s\ {\rm odd}}w_j(s),\qquad
 \widetilde T_{j,i}=p^{d_i(j)}T_{j,i}.
\]
Multiplying the degree-\(r\) identity in \eqref{eq:tag-p4} by
\(p^{d_r(j)+hr(N_j-r)}\) expresses the input operator on
every vector in the image of the chain as
\[
 Z_{j,v}K_{[j]}=p^{P_{j,v}}K_{[j]}Y_v,\qquad
 Z_{j,v}=\sum_{l=0}^r(-1)^l c_l U_j^{-(r-l)}\widetilde T_{j,r-l},
 \tag{P5}\label{eq:tag-p5}
\]
\[
 \begin{split}
 c_l&=p^{d_r(j)-d_{r-l}(j)+hl(N_j-2r+l)}h_l(\xi),\\
 P_{j,v}&=\sum_{s\ {\rm odd}}\!
      \left(rw_j(s)+\sum_{a=k-r+1}^k\lambda_{s,a}\right)
                +hrj(t+1).
 \end{split}
\]
Indeed \(Y_v=p^{A_r}T_{0,r}\), where \(A_r\) is the even head
sum minus the odd tail sum of length \(r\). Subtracting this
and the initial \(hr(k-r)\) from the multiplying exponent gives
\(P_{j,v}\). An upper bound for \(P_{j,v}\) therefore gives the required lower bound on the valuation.

\proofstep{Bound the coefficients, including the zero-index term.}

Write \(\alpha_s=\lambda_{s,1}-\lambda_{s,2}\) at a top and
zero elsewhere, and \(A_+=\sum_{s\ {\rm even}}\alpha_s\).
The weight construction has \(w_1(s)\ge-\lambda_{s,k}\)
and \(w_a(s)\) nondecreasing for \(1\le a\le m\). Hence
\[
 v_p(\xi_{ab})\ge \sum_{s\ {\rm even}}w_1(s)
                   -\sum_{s\ {\rm odd}}w_j(s)-ht\quad(a\le j).
\]
Every monomial of \(h_l\) uses exactly \(l\) entries. Substitution
in \eqref{eq:tag-p5} yields
\[
 v_p(c_l)\ge
 \sum_{s\ {\rm even}}\sum_{a=r-l+1}^{r}
                   (\lambda_{s,a}+w_1(s))
             +hl(N_j-2r+l-t)\ \ge0 .
 \tag{P6}\label{eq:tag-p6}
\]
Here \(j\ge1\) and \(k\ge2r\). At \(l=r\) the sum includes
\(\lambda_{s,1}+w_1(s)\ge\lambda_{s,1}-\lambda_{s,2}\)
for every even top. Thus \(v_p(c_r)\ge A_+\).
This estimate pays for the absence of a positive-index Hecke
operator in that summand.

\end{proof}

\paragraph{Choosing the stages on a single column.}
We spell out the remaining entry hypothesis for the use of \eqref{eq:tag-p5}.
Suppose the normalized stage kernel on the tuple at \(w\), on
integral input, has the following simultaneous bounds. Choose any
set \(B\subseteq V\) of distinct active orbits at stage \(j\), where
\(V\) is the image of the first-embedding representatives, and an
index \(0\le i_v\le r\) at each of its members. On applying \(\prod_{v\in B}
\widetilde T_{j,i_v}\) at the target, the lower valuation is
\[
 L_j+\sum_{v\notin B} b_{jv}
    +\sum_{v\in B}\big(G_{jv}-{\bf1}_{i_v=0}A_{+,v}\big).
 \tag{P7}\label{eq:tag-p7}
\]
The sums over \(v\) in \eqref{eq:tag-p7} and \eqref{eq:tag-p8} run over \(V\).

In the
range \(\log_p(2a)\ge1\), one uses
\(b_{jv}=\max(0,g^0_{jv})\) and \(G_{jv}=g^U_{jv}\),
provided the numerical lower bounds lie below the local base
and positive-index bounds and the numerical upper bound dominates
\eqref{eq:tag-p5}. In the other range both required improvements over zero
are nonpositive, so the baseline zero suffices on \(V\).
At the places outside \(V\) let \(F'_{jv}\) be the common
one-place lower bound, and let \(d_{j-1}\) bound the
stabilizer denominator in the transition. The term in \eqref{eq:tag-p7} is

\[
  L_j=\sum_{v\notin V}F'_{jv}-d_{j-1}.
\]
In the application below the cutoff,
\(F'_{jv}=F_{\rm orb}-\mathrm{err}_F\) on the second-embedding
orbit, which can have either sign. Above the cutoff it is zero.
Changing the representative rotates the letters of the cycle,
preserving the forced sum and the number of tops in the error.
Thus one may aggregate these entries by their cycles when
checking the required common lower bound against 
\[
  \sum_j L_j.
\]
 In this hypothesis a lower
bound means membership in the integral target multiplied by the
uniformizer to the ceiling of the bound times the ramification
index. The entry hypothesis involves a single lift.

It need only
be proved on basis columns and on one incident source and target
at a time, before summing their path expansion with the
stabilizer denominator.
For basis elements \(e_i\) at the source and a fixed target test \(\ell\), extension to an integral column uses
\[
  v_w\!\left(\ell\!\left(A\sum_i c_i e_i\right)\right)\ge\min_i v_w\bigl(\ell(Ae_i)\bigr)\qquad(c_i\in{\cal O}_w).
\]
 For instance if the entries expand into
an integral series in the separate local parameters times a
product of one factor per orbit, expand the integral series into
monomials. Each jet of given order uses only finitely many
monomials, each of which splits over the orbits. Individual
bounds valid for every integral monomial then add over the
orbits by the ultrametric inequality. This explains the need
for the entry bounds in frames adapted simultaneously to the
chosen neighbors.

By \eqref{eq:tag-p6}, expanding \(\prod_{v\in B}Z_{j,v}\) in \eqref{eq:tag-p7} gives the
same bound without any \(A_+\)-loss. The factors \(U_j\) are units.
For fixed \(J\) let \(Y_J=\prod_{v\in J}Y_v\).
\begin{proposition}[Valuation of a period after choosing its stages]\label{prop:period-stage-prices}
On the selected orbit set \(J\), with the local period identities
\eqref{eq:tag-p1} and \eqref{eq:tag-p7} and the action of each \(Y_v\) on the integral input lattice as in \eqref{eq:tag-c1}, write \(b\) for the transferred period of the marked column. Adding the local valuation bounds in \eqref{eq:tag-p8} gives
\[
 v_p(b)\ge\sum_{a=1}^j\left(L_a+\sum_{v\notin J}b_{av}+\sum_{v\in J}\max\{b_{av},G_{av}-P_{a,v}\}\right).
\]
 At an orbit outside the local condition only the baseline term \(b_{av}\) enters.
\end{proposition}

\begin{proof}
We prove by induction that, for every integral \(f\), the valuation of \(K_{[j]}Y_J^j f\) is at least
\[
 \sum_{a=1}^j\left(
 L_a+\sum_{v\notin J}b_{av}
     +\sum_{v\in J}\max(b_{av},G_{av}-P_{a,v})\right).
 \tag{P8}\label{eq:tag-p8}
\]
The case \(j=0\) is integrality. At step \(j\) take
\(B=\{v\in J:G_{jv}-P_{j,v}> b_{jv}\}\).
The operators at distinct orbits commute on the input, so

\[
  Y_J^j f=Y_B Y_J^{j-1}(Y_{J\setminus B}f).
\]
 The last input
in parentheses is integral by \eqref{eq:tag-c1} and the integral Weyl action.
Move \(Y_B\) through the whole chain using \eqref{eq:tag-p5}, losing
\(\sum_{v\in B}P_{j,v}\). Apply the induction hypothesis
through stage \(j-1\) and \eqref{eq:tag-p7} with its coefficient correction
at stage \(j\). Linearity adds the two valuation bounds since

\[
  \lceil x\rceil+\lceil y\rceil\ge\lceil x+y\rceil.
\]
This proves \eqref{eq:tag-p8} for the full range of \(j\).
For \(v\in I\setminus J\) the scalar \(p^{S_v}\) in the original
generator pays for the sum of the omitted positive differences
over \emph{all} stages by the definition of \(S_v\).
Taking the integral span and the trace-zero intersection
preserves this lower bound. Thus for a period with \eqref{eq:tag-p1} and \eqref{eq:tag-p7},
the column in \eqref{eq:tag-c5} realizes the full gain on the selected orbit at the cost of
the single-stage allowances in \eqref{eq:tag-p7}. At an orbit without the local condition the
base bound alone enters. The bound is uniform in the number of
monomials, words and active sets. Proposition~\ref{prop:incident-jet-indices}
proves the individual coefficient estimates. Sections~\ref{subsec:incident-sections}
and~\ref{subsec:second-period} identify the selected theta
readings and the contributions of the other cycles at a prime with a selected orbit. The bound at the assigned prime has its own marked incidence
identity \eqref{eq:tag-w2}. \label{cor:hecke-period-1}
\end{proof}

